% Paper 6 (GRPO-Registry) discussion + limitations + conclusion
\section{Discussion and Limitations}
\label{sec:p6-discussion}
Five limitations bound what this resource can claim. First, \emph{seed
breadth}: all 28 stack entries derive from one research program's
artifacts. They exercise every schema feature (open/managed, implemented/
surrogate/absent/unknown, fixed/adaptive $G$, dry-run vs.\ completed), but the
registry's value at community scale is a hypothesis until external entries
arrive. Second, \emph{evidence scale}: the open-trainer arms are toy-scale
Colab runs; we use them as existence proofs that the schema can capture a
fully-audited stack, and as directional corroboration of the manifest-level
DAPO diff---not as effect sizes. The larger-scale observations (the
under-specified same-label swing and the 368-run per-model ZVF table) come
from the program audits reported in the companion papers and are themselves
confounded (checkpoint and harness for the former; model and configuration
for the latter). Third, \emph{self-report}: for closed
stacks the registry records what the operator exposes; the null semantics and
R5 rung make unauditability visible, but they cannot manufacture the missing
facts---an entry can be honest and still thin. Fourth, \emph{badge
semantics}: the badge measures reporting coverage with equal item weights; it
is not a correctness proof, not a quality score, and alternative weightings
would reorder the middle of Table~\ref{tab:p6-entries}\ifdefined\SUBMISSIONBUILD{} (Appendix~\ref{app:p6s-entries})\fi{} (the endpoints are
robust). We ship the scoring function precisely so disagreements about
weights can be run, not argued. Fifth, \emph{curation}: some seed entries
carry errors flagged here but not yet corrected in the data.
\ifdefined\SUBMISSIONBUILD\else
\paragraph{Vehicle packaging.}
In the program's consolidation plan this resource is the appendix/linked
artifact of the MIN-REPORT-RL position paper when co-submitted, and a
standalone position-artifact paper when submitted alone. Either packaging is
valid; neither requires new training runs.
\fi

\section{Conclusion}
\label{sec:p6-conclusion}
GRPO variants are defined by small deltas that current practice scatters
across appendices, defaults, and closed launchers. Our program's audits show
why that scattering matters for interpretation: two runs sharing a label and
a short hyperparameter table differed by $17\times$ in final reward, and the
difference traces to an unreported checkpoint switch plus reward, prompt,
decoding, batch, data, and loss-implementation differences.
\textsc{GRPO-Registry} is the smallest artifact we could build that makes the
problem tractable: a schema that
carries the seven-field run-manifest portion of the eight-item reporting
standard and per-component
variant-delta status, 28 seed stack entries derived from real artifacts, and a
CLI whose \texttt{stackdiff} lists those confounds field by field, with
severity, coverage, and an underspecified flag. The registry does not make
stacks comparable; it makes their incomparability precise, and that is the
prerequisite for every fair comparison the field wants to run. Schema,
entries, and tooling are released with the benchmark; new entries arrive by
pull request (no trainer-side emitter is shipped yet).
